% 性能约束
\section{性能约束：排列模式、负载包络与无阻塞潜能}\label{sec:ov-performance}

第一类约束回答：给定端口带宽 $B$，网络在最坏流量模式下是否仍具备无阻塞的\textbf{潜能}？

\subsection{输入：排列模式集合 $\mathcal{R}$}

最坏流量模式由 Birkhoff--von Neumann 定理~\cite{birkhoff1946tres}给出：
双随机流量矩阵多面体的顶点恰为排列矩阵，线性负载在最坏情形下必在排列顶点取到。
因此我们预先固定一组排列需求矩阵 $\{\mathbf{D}^{(r)}\}_{r\in\mathcal{R}}$ 作为\textbf{外生输入}——
$\mathbf{D}$ 不是 LP 的决策变量，网络不能自选流量。
（把 $|\mathcal{R}|$ 从 $n!$ 压到可控规模的群论归约只是降低 LP 规模的实现技巧，见附录 A；
它不改变约束形式，也非本文核心贡献。）

\subsection{负载包络}

对每个模式 $r$，分流变量 $\mathbf{f}^{(r)}$ 满足流守恒，得到模式负载 $\mathbf{L}^{(r)}$；
共享包络 $\mathbf{L}$ 取所有模式的最大值：

\begin{equation}\label{eq:route-flow}
    \forall r \in \mathcal{R}:\quad
    \sum_k f_{ij}^{k,(r)} = D_{ij}^{(r)},\qquad
    \mathbf{L}^{(r)} = \mathcal{P}_r \mathbf{f}^{(r)},\qquad
    \mathbf{L} \;\ge\; \mathbf{L}^{(r)}
\end{equation}

其中 $\mathcal{P}_r$ 为模式 $r$ 的路径 incidence。包络意味着按最坏情形配置物理资源：
所有模式同时以包络负载运行，得到保守的资源与温度上界。
为压出最省资源的包络下界，单独求解性能侧时用目标 $\min\sum_e L_e$：

\begin{equation}\label{eq:valiant-lp}
    \min_{\{\mathbf{f}^{(r)}\}, \mathbf{L}} \; \sum_e L_e
    \quad\text{s.t.}\quad
    \sum_k f_{ij}^{k,(r)} = D_{ij}^{(r)},\;
    \mathbf{L}^{(r)} = \mathcal{P}_r \mathbf{f}^{(r)},\;
    \mathbf{L} \ge \mathbf{L}^{(r)} \;\;\forall r \in \mathcal{R}
\end{equation}

该 LP 只涉及拓扑与路由、与物理参数无关，输出``期待''侧的需求下界 $\mathbf{L}$。
$\mathbf{L}$ 是整张 LP 的枢纽：一端接收 $\mathcal{R}$ 个排列模式的路由需求，
另一端经 $\boldsymbol{\ell} = B\,\mathbf{S}_{\text{bw}}^{-1}\mathbf{L}$ 转化为物理代价。
$\mathbf{L}$ 往上走 $\to$ $\boldsymbol{\ell}$ 涨 $\to$ bump、布线、热约束收紧。

\textbf{关键观察}：性能约束的形式始终是关于 $\mathbf{L}$ 的线性等式/不等式；
$\mathbf{D}^{(r)}$ 外生固定，无阻塞语义由包络 $\mathbf{L}\ge\mathbf{L}^{(r)}$ 承担。
